% Context from chapters/denoising.tex; for mathematical reference, not compilation.
\section{Linear Denoising by Filtering}
\label{sec-linear-denoising}

                                                     
\subsection{Translation Invariant Estimators}

A linear estimator $\Dd(Y) = \tilde f$ of $f_0$ preserves addition, $\Dd(f+g) = \Dd(f)+\Dd(g)$, and scalar multiplication.
A translation-invariant estimator commutes with translations: $\Dd(f_\tau) = \Dd(f)_\tau$, where $f_\tau(t) = f(t-\tau)$. We assume periodic boundary conditions, in either one or two dimensions.
 
An estimator with both properties is a convolution filter:
\eq{
	\Dd(f) = f \star h
}
where $h\in\RR^N$ is a filter. Denoising filters are usually low-pass. To preserve constant signals, one imposes
\eq{
	\sum_n h_n = \hat h_0 = 1
}
where $\hat h$ is the discrete Fourier transform. 

Figure \ref{fig-filtering-denoised} illustrates denoising with a low-pass filter.

\myfigure{
\tabdeux{
\image{denoising}{.45}{filtering-noisy}&
\image{denoising}{.45}{filtering-noisy-fourier}\\
$f$ & $\log|\hat f|$ \\
\image{denoising}{.45}{filtering-filter}&
\image{denoising}{.45}{filtering-filter-fourier}\\
$h$ & $\log|\hat h|$ \\
\image{denoising}{.45}{filtering-denoised}&
\image{denoising}{.45}{filtering-denoised-fourier}\\
$f \star h$ & $\log|\hat f\hat h|=\log|\hat f|+\log|\hat h|$
}
}{ 
	Filtering in the spatial domain (left) and the Fourier domain (right). 
}{fig-filtering-denoised}

The filtering strength is usually controlled by the width $s$ of $h$. A typical example is the (discretized) Gaussian filter
\eql{\label{eq-gaussian-filters}
	\foralls -N/2 < i \leq N/2, \quad h_{s,i} = \frac{1}{Z_s} \exp\pa{ -\frac{ i^2 }{ 2s^2} }
}
where $Z_s$ normalizes the filter so that $\sum_i h_{s,i}=1$, preserving constant signals. Figure \ref{fig-filtering-denoised} shows the effect of Gaussian filtering in the spatial and Fourier domains.

     
          
                                          
                                           
   
                                                                    
                  
          
                                        
                                           
   
                                                                  
                  
    

Figure \ref{fig-filtering-progression} shows the effect of increasing the filter width $s$. Linear filtering works well on smooth signals, but removing more noise also blurs edges and other singularities.

\myfigure{
\image{denoising}{.7}{filtering-progression-1d}
\image{denoising}{.2}{filtering-progression-2d}
}{ 
	Denoising using a filter of increasing width $s$.  
}{fig-filtering-progression}

                                                     
\subsection{Optimal Filter Selection and Bias-Variance Tradeoff}
\label{subsec-optimal-selection-filter}

The best filter depends on both the unknown signal $f_0$ and the noise level $\si$. Rather than optimize over all filters, we can tune a width parameter $s$ within a family such as the Gaussian filters in~\eqref{eq-gaussian-filters}.

The triangle inequality separates the smoothing error from the remaining noise:
\eq{
	\norm{\tilde f - f_0} \leq \norm{ h_s \star f_0 - f_0 } + \norm{ h_s \star w }
}
The filter width $s$ balances noise removal against excessive smoothing of singularities. A wider filter typically reduces $\norm{h_s \star w}$ while increasing $\norm{h_s\star f_0-f_0}$.

Figur